% GENERATED FILE. Edit canonical lessons, then run npm run generate:books.
% canonical-source-hash: fnv1a64:ca434f5bf8baca73
% canonical-lessons: KA-C435-bart, KA-C435-asana, KA-C435-kaitappu, KA-C435-inuku, KA-C435-dara

\chapter{Berths and Missed Trains}
\label{ch:ka-berths-and-missed-trains}
\hlchaptermodality{435}

\begin{chapteropening}
\textbf{By the end of this chapter:} \emph{I can say a berth, a formal seat, to slip away, to peep and a fare.}

Complete the last lesson of chapter 435: I can say a berth, a formal seat, to slip away, to peep and a fare.
\end{chapteropening}

\section[\texorpdfstring{\kn{ಬರ್ತ್} (bart)}{ಬರ್ತ್ (bart)}]{\kn{ಬರ್ತ್} (bart) — a berth}

\label{lesson:KA-C435-bart}



\begin{quote}
Before the new one: say the Kannada for a halt, then the Kannada for to enquire.
\end{quote}

\subsection*{You'll want to know: \kn{ಬರ್ತ್}}
\textbf{\kn{ಬರ್ತ್}} — \emph{bart} — \textquotedblleft{}a berth\textquotedblright{}.

\begin{quote}
\kn{ನನಗೆ} \kn{ಮೇಲಿನ} \kn{ಬರ್ತ್} \kn{ಸಿಕ್ಕಿತು}.
\end{quote}

\begin{quote}
\emph{nanage mēlina bart sikkitu.}
\end{quote}

\begin{quote}
I got the upper berth.
\end{quote}

\kn{ಮೇಲಿನ}, \kn{ನಡುವಿನ}, \kn{ಕೆಳಗಿನ} \kn{ಬರ್ತ್}: the upper, middle and lower berth.

\begin{cousinweb}
\textbf{\kn{ಬರ್ತ್}} (\emph{bart}) is a loan from English \emph{berth}, a sleeping place on a train or a ship.
\end{cousinweb}

\begin{grammarlens}[title={What you've built}]
One more word for getting around.
\end{grammarlens}

\subsection*{Your turn}
\begin{itemize}
\raggedright
  \item \emph{Say it:} \emph{bart}
  \item \emph{Say it:} \emph{bart}, once more
  \item \emph{Say it:} \emph{nanage mēlina bart sikkitu.} — I got the upper berth.
  \item \emph{Recall:} say the Kannada for a halt, then the Kannada for to enquire, then say \emph{bart} again
\end{itemize}

\begin{tcolorbox}[breakable,colback=teal!4,colframe=teal!35!black,title={Before you move on}]
What does \kn{ಬರ್ತ್} mean? (\textquotedblleft{}A berth\textquotedblright{}.) Say it once more.
\end{tcolorbox}
\section[\texorpdfstring{\kn{ಆಸನ} (āsana)}{ಆಸನ (āsana)}]{\kn{ಆಸನ} (āsana) — a formal seat; a posture}

\label{lesson:KA-C435-asana}



\begin{quote}
Before the new one: say the Kannada for to enquire, then the Kannada for a berth.
\end{quote}

\subsection*{You'll want to know: \kn{ಆಸನ}}
\textbf{\kn{ಆಸನ}} — \emph{āsana} — \textquotedblleft{}a formal seat; a posture\textquotedblright{}.

\begin{quote}
\kn{ಬಸ್ಸಿನಲ್ಲಿ} \kn{ಹಿರಿಯರಿಗೆ} \kn{ಬೇರೆ} \kn{ಆಸನಗಳು} \kn{ಇವೆ}.
\end{quote}

\begin{quote}
\emph{bassinalli hiriyarige bēre āsanagaḷu ive.}
\end{quote}

\begin{quote}
On the bus there are separate seats for elders.
\end{quote}

Plural: \kn{ಆಸನಗಳು}, \emph{āsanagaḷu}.

The formal word of signs and tickets, where everyday speech says \kn{ಸೀಟು}; in yoga an \kn{ಆಸನ} is a posture.

\begin{cousinweb}
\textbf{\kn{ಆಸನ}} (\emph{āsana}) is borrowed from Sanskrit \emph{āsana}, \textquotedblleft{}sitting, a seat\textquotedblright{} (CDIAL 1484), from the root \emph{ās}, \textquotedblleft{}to sit\textquotedblright{}.
\end{cousinweb}

\begin{grammarlens}[title={What you've built}]
One more word for getting around.
\end{grammarlens}

\subsection*{Your turn}
\begin{itemize}
\raggedright
  \item \emph{Say it:} \emph{āsana}
  \item \emph{Say it:} \emph{āsana}, once more
  \item \emph{Say it:} \emph{bassinalli hiriyarige bēre āsanagaḷu ive.} — On the bus there are separate seats for elders.
  \item \emph{Recall:} say the Kannada for to enquire, then the Kannada for a berth, then say \emph{āsana} again
\end{itemize}

\begin{tcolorbox}[breakable,colback=teal!4,colframe=teal!35!black,title={Before you move on}]
What does \kn{ಆಸನ} mean? (\textquotedblleft{}A formal seat; a posture\textquotedblright{}.) Say it once more.
\end{tcolorbox}
\section[\texorpdfstring{\kn{ಕೈತಪ್ಪು} (kaitappu)}{ಕೈತಪ್ಪು (kaitappu)}]{\kn{ಕೈತಪ್ಪು} (kaitappu) — to slip away; to be missed}

\label{lesson:KA-C435-kaitappu}



\begin{quote}
Before the new one: say the Kannada for a berth, then the Kannada for a formal seat.
\end{quote}

\subsection*{You'll want to know: \kn{ಕೈತಪ್ಪು}}
\textbf{\kn{ಕೈತಪ್ಪು}} — \emph{kaitappu} — \textquotedblleft{}to slip away; to be missed\textquotedblright{}.

\begin{quote}
\kn{ಎರಡು} \kn{ನಿಮಿಷ} \kn{ತಡವಾದ್ದರಿಂದ} \kn{ರೈಲು} \kn{ಕೈತಪ್ಪಿತು}.
\end{quote}

\begin{quote}
\emph{eraḍu nimiṣa taḍavāddarinda railu kaitappitu.}
\end{quote}

\begin{quote}
Being two minutes late, I missed the train.
\end{quote}

Past (it …): \kn{ಕೈತಪ್ಪಿತು}, \emph{kaitappitu}.

The thing missed is the subject: \kn{ಅವಕಾಶ} \kn{ಕೈತಪ್ಪಿತು}, the chance slipped away.

\begin{cousinweb}
\textbf{\kn{ಕೈತಪ್ಪು}} (\emph{kaitappu}) is a native Kannada compound: \emph{kai}, \textquotedblleft{}hand\textquotedblright{}, and \kn{ತಪ್ಪು}, \textquotedblleft{}to slip, to miss\textquotedblright{} (DEDR 3071): to slip out of the hand.
\end{cousinweb}

\begin{grammarlens}[title={What you've built}]
One more word for getting around.
\end{grammarlens}

\subsection*{Your turn}
\begin{itemize}
\raggedright
  \item \emph{Say it:} \emph{kaitappu}
  \item \emph{Say it:} \emph{kaitappu}, once more
  \item \emph{Say it:} \emph{eraḍu nimiṣa taḍavāddarinda railu kaitappitu.} — Being two minutes late, I missed the train.
  \item \emph{Recall:} say the Kannada for a berth, then the Kannada for a formal seat, then say \emph{kaitappu} again
\end{itemize}

\begin{tcolorbox}[breakable,colback=teal!4,colframe=teal!35!black,title={Before you move on}]
What does \kn{ಕೈತಪ್ಪು} mean? (\textquotedblleft{}To slip away; to be missed\textquotedblright{}.) Say it once more.
\end{tcolorbox}
\section[\texorpdfstring{\kn{ಇಣುಕು} (iṇuku)}{ಇಣುಕು (iṇuku)}]{\kn{ಇಣುಕು} (iṇuku) — to peep; to peer}

\label{lesson:KA-C435-inuku}



\begin{quote}
Before the new one: say the Kannada for a formal seat, then the Kannada for to slip away.
\end{quote}

\subsection*{You'll want to know: \kn{ಇಣುಕು}}
\textbf{\kn{ಇಣುಕು}} — \emph{iṇuku} — \textquotedblleft{}to peep; to peer\textquotedblright{}.

\begin{quote}
\kn{ಮಗು} \kn{ಕಿಟಕಿಯಿಂದ} \kn{ಹೊರಗೆ} \kn{ಇಣುಕಿತು}.
\end{quote}

\begin{quote}
\emph{magu kiṭakiyinda horage iṇukitu.}
\end{quote}

\begin{quote}
The child peeped out of the window.
\end{quote}

Past (I …): \kn{ಇಣುಕಿದೆ}, \emph{iṇukide}.

\kn{ಇಣುಕಿ} \kn{ನೋಡು}, to peep and look, is the common pair.

\begin{cousinweb}
\textbf{\kn{ಇಣುಕು}} (\emph{iṇuku}) is a Kannada verb whose origin is not settled; DEDR shows no clear source for it.
\end{cousinweb}

\begin{grammarlens}[title={What you've built}]
One more word for getting around.
\end{grammarlens}

\subsection*{Your turn}
\begin{itemize}
\raggedright
  \item \emph{Say it:} \emph{iṇuku}
  \item \emph{Say it:} \emph{iṇuku}, once more
  \item \emph{Say it:} \emph{magu kiṭakiyinda horage iṇukitu.} — The child peeped out of the window.
  \item \emph{Recall:} say the Kannada for a formal seat, then the Kannada for to slip away, then say \emph{iṇuku} again
\end{itemize}

\begin{tcolorbox}[breakable,colback=teal!4,colframe=teal!35!black,title={Before you move on}]
What does \kn{ಇಣುಕು} mean? (\textquotedblleft{}To peep; to peer\textquotedblright{}.) Say it once more.
\end{tcolorbox}
\section[\texorpdfstring{\kn{ದರ} (dara)}{ದರ (dara)}]{\kn{ದರ} (dara) — a fare; a rate}

\label{lesson:KA-C435-dara}



\begin{quote}
Before the new one: say the Kannada for to slip away, then the Kannada for to peep.
\end{quote}

\subsection*{You'll want to know: \kn{ದರ}}
\textbf{\kn{ದರ}} — \emph{dara} — \textquotedblleft{}a fare; a rate\textquotedblright{}.

\begin{quote}
\kn{ಈ} \kn{ವರ್ಷ} \kn{ಬಸ್} \kn{ದರ} \kn{ಹೆಚ್ಚಾಯಿತು}.
\end{quote}

\begin{quote}
\emph{ī varṣa bas dara heccāyitu.}
\end{quote}

\begin{quote}
This year the bus fare went up.
\end{quote}

\kn{ಪ್ರಯಾಣ} \kn{ದರ} is a fare; \kn{ಬಡ್ಡಿ} \kn{ದರ} is a rate of interest.

\begin{cousinweb}
\textbf{\kn{ದರ}} (\emph{dara}) is usually traced to Hindustani \emph{dar}, \textquotedblleft{}rate, price\textquotedblright{}; its earlier source is not settled.
\end{cousinweb}

\begin{grammarlens}[title={What you've built}]
One more word for getting around.
\end{grammarlens}

\subsection*{Your turn}
\begin{itemize}
\raggedright
  \item \emph{Say it:} \emph{dara}
  \item \emph{Say it:} \emph{dara}, once more
  \item \emph{Say it:} \emph{ī varṣa bas dara heccāyitu.} — This year the bus fare went up.
  \item \emph{Recall:} say the Kannada for to slip away, then the Kannada for to peep, then say \emph{dara} again
\end{itemize}

\begin{tcolorbox}[breakable,colback=teal!4,colframe=teal!35!black,title={Before you move on}]
What does \kn{ದರ} mean? (\textquotedblleft{}A fare; a rate\textquotedblright{}.) Say it once more.
\end{tcolorbox}
